\section{Related Work}\label{sec:rel}

\toolname occupies a specific niche in the binary-analysis landscape: it associates \textit{known} input-format labels with a binary's basic blocks using differential QEMU translation logs.
We position it against three lines of prior work that are easily confused with each other but solve materially different problems.

\subsection{Trace-Diffing Tools}\label{sec:rel:tenet}

Tenet~\cite{gaasedelen2021tenet} and Lighthouse~\cite{gaasedelen2017lighthouse} are disassembler plugins that highlight the basic blocks differing between execution traces of the same binary inside IDA~Pro or Binary~Ninja.
They are widely used in practice for tasks like identifying why one input crashes while another does not.

These tools solve a degenerate case of \toolname's problem: $N{=}2$ inputs, no input-side annotations.
\toolname generalizes the differential approach in two ways.
First, we systematically generate large corpora of structurally related inputs from a grammar template (Section~\ref{sec:design:corpus}) rather than asking the user to supply two interesting inputs by hand.
Second, we attach semantic labels to the input structure and propagate those labels to trace-derived record pairs, producing a format-aware annotation rather than a binary ``differs / does not differ'' verdict.

\subsection{Input-Region Categorization}\label{sec:rel:polyfile}

PolyFile~\cite{sultanik2019polyfile} from Trail of Bits identifies which byte ranges of a file belong to which substructure of a known file format.
PolyTracker~\cite{sultanik2024polytracker} extends this with whole-input dynamic information flow tracking, using LLVM instrumentation to build a directed acyclic graph that maps each input byte through the program's execution.

\toolname and these tools both operate on inputs whose formats are known, but they differ in what they annotate.
PolyFile produces an annotation of the \textit{input}.
PolyTracker produces a data-flow mapping between input bytes and program operations.
\toolname produces candidate annotations for binary locations represented in its trace-derived graph.
The three tools are therefore composable: PolyFile can produce the labeled corpus that \toolname consumes, and \toolname's output can be cross-referenced with PolyTracker's data-flow graph for deeper understanding.

A methodological distinction separates \toolname from PolyTracker: PolyTracker uses LLVM-level instrumentation, whereas \toolname observes an executable through \QEMU user-mode emulation.
The resulting tradeoff is between richer byte-level data flow and a lighter-weight workflow that can be applied when a runnable binary and compatible emulation environment are available.

\subsection{Grammar Mining and Format Recovery}\label{sec:rel:grammar}

A substantial body of academic work addresses the inverse problem: given only a binary, recover the grammar of the inputs it accepts.
Tupni~\cite{cui2008tupni} performs automatic reverse engineering of input formats using dynamic analysis.
Polyglot~\cite{caballero2007polyglot} extracts protocol message formats via dynamic binary analysis.
AUTOGRAM~\cite{hoschele2016autogram} infers grammars from observed input-output behavior using dynamic taint analysis.
More recent work by Gopinath et~al.~\cite{gopinath2020mining} mines context-free input grammars from dynamic control flow.
Prospex~\cite{comparetti2009prospex} recovers protocol state machines from network traces.

These approaches answer the question \textit{``what does this binary eat?''}
\toolname answers a different question: \textit{``given that this binary eats X, where in the code is X handled?''}
The two problems are sequential rather than competing.
An analyst would use a grammar-recovery tool to obtain a starting grammar, then use \toolname to map that grammar onto the code.
Conflating the two would be a category error.

A second methodological distinction concerns observation cost and evidence.
Grammar-recovery systems span compiler instrumentation, dynamic taint analysis, and binary-level execution analysis; several, including Tupni and Polyglot, operate directly on binaries.
\toolname does not recover an unknown grammar or track byte-level data flow.
It instead assumes a known template and uses lightweight \QEMU translation-log differences, trading semantic depth for easier deployment on compatible runnable binaries.

\subsection{Differential Testing for Parsers}\label{sec:rel:difftesting}

The broader idea of detecting parser bugs through differential observation has been explored in the HTTP/1.1 ecosystem~\cite{http-smuggling} and in static differential analysis of protocol parsers~\cite{zheng2024pardiff}.
These tools surface discrepancies in protocol semantics between parser implementations.
\toolname surfaces the structural correspondence between an agreed input grammar and the code that realizes it.
They are complementary: a differential-testing tool can find that two parsers disagree, then \toolname can localize where in each parser the disagreement lives.

% TODO: Revisit the differential-parser citation set before publication.
